% ============================================================================
% ESQUEMA (banca): anatomia de uma chamada de rotulagem em lote ao oráculo
% LLM — prefixo estático cacheável + itens numerados -> resposta como vetor
% indexado. Inserido no APÊNDICE E (a5-prompts/texto.tex) por ordem do
% autor (2026-08-23), dentro de figure:
%   % ============================================================================
% ESQUEMA (banca): anatomia de uma chamada de rotulagem em lote ao oráculo
% LLM — prefixo estático cacheável + itens numerados -> resposta como vetor
% indexado. Inserido no APÊNDICE E (a5-prompts/texto.tex) por ordem do
% autor (2026-08-23), dentro de figure:
%   % ============================================================================
% ESQUEMA (banca): anatomia de uma chamada de rotulagem em lote ao oráculo
% LLM — prefixo estático cacheável + itens numerados -> resposta como vetor
% indexado. Inserido no APÊNDICE E (a5-prompts/texto.tex) por ordem do
% autor (2026-08-23), dentro de figure:
%   % ============================================================================
% ESQUEMA (banca): anatomia de uma chamada de rotulagem em lote ao oráculo
% LLM — prefixo estático cacheável + itens numerados -> resposta como vetor
% indexado. Inserido no APÊNDICE E (a5-prompts/texto.tex) por ordem do
% autor (2026-08-23), dentro de figure:
%   \input{3-metodo/esquemas-propostos/esq-prompt-anatomia}
% AUTOCONTIDO: sem \ref, sem códigos de experimento; o único número (621)
% é o do CategorySchema, definido na tese e usado no próprio apêndice.
% Uso standalone (pré-visualização):
%   \documentclass[tikz,border=6pt,12pt]{standalone}
%   \usetikzlibrary{arrows.meta, positioning}
%   \begin{document}\input{esq-prompt-anatomia.tex}\end{document}
% GEOMETRIA CALIBRADA PARA CORPO 12 (ppginf/book 12pt, textwidth 16cm).
% Requer apenas arrows.meta e positioning (já em packages.tex). Legível em
% P&B: prefixo cacheável = cinza claro; resposta = cinza médio;
% contorno tracejado = a chamada.
% ============================================================================
\begin{tikzpicture}[
    x=1cm, y=1cm,
    caixa/.style={draw, rounded corners=2pt, align=center, text width=56mm,
                  minimum height=12mm, inner sep=3pt, font=\footnotesize},
    pref/.style={caixa, fill=gray!8},
    resp/.style={caixa, fill=gray!14},
    seta/.style={-{Stealth[length=2.6mm]}, thick},
    rot/.style={font=\scriptsize, align=center, inner sep=2pt}
]
    % ---------------- a chamada (contorno tracejado) -----------------------
    \node[pref] (prefixo) at (0,0)
        {\textbf{prefixo estático}, idêntico\\entre chamadas: contexto\\
         da tarefa $+$ esquema de saída};
    \node[caixa] (itens) at (0,-2.6)
        {\textbf{itens do lote}, numerados\\(1 a $n$): as descrições\\
         a rotular};
    \draw[dashed, rounded corners=3pt] (-3.25,1.15) rectangle (3.25,-3.75);
    \node[rot, anchor=south west] at (-3.2,1.25) {uma chamada ao provedor};

    % ---------------- oráculo e resposta -----------------------------------
    \node[caixa, text width=28mm] (llm) at (5.6,-1.3)
        {oráculo LLM\\(provedor)};
    \node[resp, text width=48mm] (vet) at (10.2,-1.3)
        {\textbf{resposta}: vetor indexado,\\um objeto por item\\
         (categoria restrita às 621;\\expansão; justificativa)};

    \draw[seta] (3.25,-1.3) -- (llm.west);
    \draw[seta] (llm.east) -- (vet.west);

    % ---------------- nota do cache ----------------------------------------
    \node[rot, anchor=west, text width=128mm, align=left] at (-3.2,-4.7)
        {o prefixo repetido habilita o \textit{cache} de prefixo dos
         provedores; o esquema do lote não fixa o número de itens,
         preservando a cacheabilidade.};
\end{tikzpicture}

% AUTOCONTIDO: sem \ref, sem códigos de experimento; o único número (621)
% é o do CategorySchema, definido na tese e usado no próprio apêndice.
% Uso standalone (pré-visualização):
%   \documentclass[tikz,border=6pt,12pt]{standalone}
%   \usetikzlibrary{arrows.meta, positioning}
%   \begin{document}% ============================================================================
% ESQUEMA (banca): anatomia de uma chamada de rotulagem em lote ao oráculo
% LLM — prefixo estático cacheável + itens numerados -> resposta como vetor
% indexado. Inserido no APÊNDICE E (a5-prompts/texto.tex) por ordem do
% autor (2026-08-23), dentro de figure:
%   \input{3-metodo/esquemas-propostos/esq-prompt-anatomia}
% AUTOCONTIDO: sem \ref, sem códigos de experimento; o único número (621)
% é o do CategorySchema, definido na tese e usado no próprio apêndice.
% Uso standalone (pré-visualização):
%   \documentclass[tikz,border=6pt,12pt]{standalone}
%   \usetikzlibrary{arrows.meta, positioning}
%   \begin{document}\input{esq-prompt-anatomia.tex}\end{document}
% GEOMETRIA CALIBRADA PARA CORPO 12 (ppginf/book 12pt, textwidth 16cm).
% Requer apenas arrows.meta e positioning (já em packages.tex). Legível em
% P&B: prefixo cacheável = cinza claro; resposta = cinza médio;
% contorno tracejado = a chamada.
% ============================================================================
\begin{tikzpicture}[
    x=1cm, y=1cm,
    caixa/.style={draw, rounded corners=2pt, align=center, text width=56mm,
                  minimum height=12mm, inner sep=3pt, font=\footnotesize},
    pref/.style={caixa, fill=gray!8},
    resp/.style={caixa, fill=gray!14},
    seta/.style={-{Stealth[length=2.6mm]}, thick},
    rot/.style={font=\scriptsize, align=center, inner sep=2pt}
]
    % ---------------- a chamada (contorno tracejado) -----------------------
    \node[pref] (prefixo) at (0,0)
        {\textbf{prefixo estático}, idêntico\\entre chamadas: contexto\\
         da tarefa $+$ esquema de saída};
    \node[caixa] (itens) at (0,-2.6)
        {\textbf{itens do lote}, numerados\\(1 a $n$): as descrições\\
         a rotular};
    \draw[dashed, rounded corners=3pt] (-3.25,1.15) rectangle (3.25,-3.75);
    \node[rot, anchor=south west] at (-3.2,1.25) {uma chamada ao provedor};

    % ---------------- oráculo e resposta -----------------------------------
    \node[caixa, text width=28mm] (llm) at (5.6,-1.3)
        {oráculo LLM\\(provedor)};
    \node[resp, text width=48mm] (vet) at (10.2,-1.3)
        {\textbf{resposta}: vetor indexado,\\um objeto por item\\
         (categoria restrita às 621;\\expansão; justificativa)};

    \draw[seta] (3.25,-1.3) -- (llm.west);
    \draw[seta] (llm.east) -- (vet.west);

    % ---------------- nota do cache ----------------------------------------
    \node[rot, anchor=west, text width=128mm, align=left] at (-3.2,-4.7)
        {o prefixo repetido habilita o \textit{cache} de prefixo dos
         provedores; o esquema do lote não fixa o número de itens,
         preservando a cacheabilidade.};
\end{tikzpicture}
\end{document}
% GEOMETRIA CALIBRADA PARA CORPO 12 (ppginf/book 12pt, textwidth 16cm).
% Requer apenas arrows.meta e positioning (já em packages.tex). Legível em
% P&B: prefixo cacheável = cinza claro; resposta = cinza médio;
% contorno tracejado = a chamada.
% ============================================================================
\begin{tikzpicture}[
    x=1cm, y=1cm,
    caixa/.style={draw, rounded corners=2pt, align=center, text width=56mm,
                  minimum height=12mm, inner sep=3pt, font=\footnotesize},
    pref/.style={caixa, fill=gray!8},
    resp/.style={caixa, fill=gray!14},
    seta/.style={-{Stealth[length=2.6mm]}, thick},
    rot/.style={font=\scriptsize, align=center, inner sep=2pt}
]
    % ---------------- a chamada (contorno tracejado) -----------------------
    \node[pref] (prefixo) at (0,0)
        {\textbf{prefixo estático}, idêntico\\entre chamadas: contexto\\
         da tarefa $+$ esquema de saída};
    \node[caixa] (itens) at (0,-2.6)
        {\textbf{itens do lote}, numerados\\(1 a $n$): as descrições\\
         a rotular};
    \draw[dashed, rounded corners=3pt] (-3.25,1.15) rectangle (3.25,-3.75);
    \node[rot, anchor=south west] at (-3.2,1.25) {uma chamada ao provedor};

    % ---------------- oráculo e resposta -----------------------------------
    \node[caixa, text width=28mm] (llm) at (5.6,-1.3)
        {oráculo LLM\\(provedor)};
    \node[resp, text width=48mm] (vet) at (10.2,-1.3)
        {\textbf{resposta}: vetor indexado,\\um objeto por item\\
         (categoria restrita às 621;\\expansão; justificativa)};

    \draw[seta] (3.25,-1.3) -- (llm.west);
    \draw[seta] (llm.east) -- (vet.west);

    % ---------------- nota do cache ----------------------------------------
    \node[rot, anchor=west, text width=128mm, align=left] at (-3.2,-4.7)
        {o prefixo repetido habilita o \textit{cache} de prefixo dos
         provedores; o esquema do lote não fixa o número de itens,
         preservando a cacheabilidade.};
\end{tikzpicture}

% AUTOCONTIDO: sem \ref, sem códigos de experimento; o único número (621)
% é o do CategorySchema, definido na tese e usado no próprio apêndice.
% Uso standalone (pré-visualização):
%   \documentclass[tikz,border=6pt,12pt]{standalone}
%   \usetikzlibrary{arrows.meta, positioning}
%   \begin{document}% ============================================================================
% ESQUEMA (banca): anatomia de uma chamada de rotulagem em lote ao oráculo
% LLM — prefixo estático cacheável + itens numerados -> resposta como vetor
% indexado. Inserido no APÊNDICE E (a5-prompts/texto.tex) por ordem do
% autor (2026-08-23), dentro de figure:
%   % ============================================================================
% ESQUEMA (banca): anatomia de uma chamada de rotulagem em lote ao oráculo
% LLM — prefixo estático cacheável + itens numerados -> resposta como vetor
% indexado. Inserido no APÊNDICE E (a5-prompts/texto.tex) por ordem do
% autor (2026-08-23), dentro de figure:
%   \input{3-metodo/esquemas-propostos/esq-prompt-anatomia}
% AUTOCONTIDO: sem \ref, sem códigos de experimento; o único número (621)
% é o do CategorySchema, definido na tese e usado no próprio apêndice.
% Uso standalone (pré-visualização):
%   \documentclass[tikz,border=6pt,12pt]{standalone}
%   \usetikzlibrary{arrows.meta, positioning}
%   \begin{document}\input{esq-prompt-anatomia.tex}\end{document}
% GEOMETRIA CALIBRADA PARA CORPO 12 (ppginf/book 12pt, textwidth 16cm).
% Requer apenas arrows.meta e positioning (já em packages.tex). Legível em
% P&B: prefixo cacheável = cinza claro; resposta = cinza médio;
% contorno tracejado = a chamada.
% ============================================================================
\begin{tikzpicture}[
    x=1cm, y=1cm,
    caixa/.style={draw, rounded corners=2pt, align=center, text width=56mm,
                  minimum height=12mm, inner sep=3pt, font=\footnotesize},
    pref/.style={caixa, fill=gray!8},
    resp/.style={caixa, fill=gray!14},
    seta/.style={-{Stealth[length=2.6mm]}, thick},
    rot/.style={font=\scriptsize, align=center, inner sep=2pt}
]
    % ---------------- a chamada (contorno tracejado) -----------------------
    \node[pref] (prefixo) at (0,0)
        {\textbf{prefixo estático}, idêntico\\entre chamadas: contexto\\
         da tarefa $+$ esquema de saída};
    \node[caixa] (itens) at (0,-2.6)
        {\textbf{itens do lote}, numerados\\(1 a $n$): as descrições\\
         a rotular};
    \draw[dashed, rounded corners=3pt] (-3.25,1.15) rectangle (3.25,-3.75);
    \node[rot, anchor=south west] at (-3.2,1.25) {uma chamada ao provedor};

    % ---------------- oráculo e resposta -----------------------------------
    \node[caixa, text width=28mm] (llm) at (5.6,-1.3)
        {oráculo LLM\\(provedor)};
    \node[resp, text width=48mm] (vet) at (10.2,-1.3)
        {\textbf{resposta}: vetor indexado,\\um objeto por item\\
         (categoria restrita às 621;\\expansão; justificativa)};

    \draw[seta] (3.25,-1.3) -- (llm.west);
    \draw[seta] (llm.east) -- (vet.west);

    % ---------------- nota do cache ----------------------------------------
    \node[rot, anchor=west, text width=128mm, align=left] at (-3.2,-4.7)
        {o prefixo repetido habilita o \textit{cache} de prefixo dos
         provedores; o esquema do lote não fixa o número de itens,
         preservando a cacheabilidade.};
\end{tikzpicture}

% AUTOCONTIDO: sem \ref, sem códigos de experimento; o único número (621)
% é o do CategorySchema, definido na tese e usado no próprio apêndice.
% Uso standalone (pré-visualização):
%   \documentclass[tikz,border=6pt,12pt]{standalone}
%   \usetikzlibrary{arrows.meta, positioning}
%   \begin{document}% ============================================================================
% ESQUEMA (banca): anatomia de uma chamada de rotulagem em lote ao oráculo
% LLM — prefixo estático cacheável + itens numerados -> resposta como vetor
% indexado. Inserido no APÊNDICE E (a5-prompts/texto.tex) por ordem do
% autor (2026-08-23), dentro de figure:
%   \input{3-metodo/esquemas-propostos/esq-prompt-anatomia}
% AUTOCONTIDO: sem \ref, sem códigos de experimento; o único número (621)
% é o do CategorySchema, definido na tese e usado no próprio apêndice.
% Uso standalone (pré-visualização):
%   \documentclass[tikz,border=6pt,12pt]{standalone}
%   \usetikzlibrary{arrows.meta, positioning}
%   \begin{document}\input{esq-prompt-anatomia.tex}\end{document}
% GEOMETRIA CALIBRADA PARA CORPO 12 (ppginf/book 12pt, textwidth 16cm).
% Requer apenas arrows.meta e positioning (já em packages.tex). Legível em
% P&B: prefixo cacheável = cinza claro; resposta = cinza médio;
% contorno tracejado = a chamada.
% ============================================================================
\begin{tikzpicture}[
    x=1cm, y=1cm,
    caixa/.style={draw, rounded corners=2pt, align=center, text width=56mm,
                  minimum height=12mm, inner sep=3pt, font=\footnotesize},
    pref/.style={caixa, fill=gray!8},
    resp/.style={caixa, fill=gray!14},
    seta/.style={-{Stealth[length=2.6mm]}, thick},
    rot/.style={font=\scriptsize, align=center, inner sep=2pt}
]
    % ---------------- a chamada (contorno tracejado) -----------------------
    \node[pref] (prefixo) at (0,0)
        {\textbf{prefixo estático}, idêntico\\entre chamadas: contexto\\
         da tarefa $+$ esquema de saída};
    \node[caixa] (itens) at (0,-2.6)
        {\textbf{itens do lote}, numerados\\(1 a $n$): as descrições\\
         a rotular};
    \draw[dashed, rounded corners=3pt] (-3.25,1.15) rectangle (3.25,-3.75);
    \node[rot, anchor=south west] at (-3.2,1.25) {uma chamada ao provedor};

    % ---------------- oráculo e resposta -----------------------------------
    \node[caixa, text width=28mm] (llm) at (5.6,-1.3)
        {oráculo LLM\\(provedor)};
    \node[resp, text width=48mm] (vet) at (10.2,-1.3)
        {\textbf{resposta}: vetor indexado,\\um objeto por item\\
         (categoria restrita às 621;\\expansão; justificativa)};

    \draw[seta] (3.25,-1.3) -- (llm.west);
    \draw[seta] (llm.east) -- (vet.west);

    % ---------------- nota do cache ----------------------------------------
    \node[rot, anchor=west, text width=128mm, align=left] at (-3.2,-4.7)
        {o prefixo repetido habilita o \textit{cache} de prefixo dos
         provedores; o esquema do lote não fixa o número de itens,
         preservando a cacheabilidade.};
\end{tikzpicture}
\end{document}
% GEOMETRIA CALIBRADA PARA CORPO 12 (ppginf/book 12pt, textwidth 16cm).
% Requer apenas arrows.meta e positioning (já em packages.tex). Legível em
% P&B: prefixo cacheável = cinza claro; resposta = cinza médio;
% contorno tracejado = a chamada.
% ============================================================================
\begin{tikzpicture}[
    x=1cm, y=1cm,
    caixa/.style={draw, rounded corners=2pt, align=center, text width=56mm,
                  minimum height=12mm, inner sep=3pt, font=\footnotesize},
    pref/.style={caixa, fill=gray!8},
    resp/.style={caixa, fill=gray!14},
    seta/.style={-{Stealth[length=2.6mm]}, thick},
    rot/.style={font=\scriptsize, align=center, inner sep=2pt}
]
    % ---------------- a chamada (contorno tracejado) -----------------------
    \node[pref] (prefixo) at (0,0)
        {\textbf{prefixo estático}, idêntico\\entre chamadas: contexto\\
         da tarefa $+$ esquema de saída};
    \node[caixa] (itens) at (0,-2.6)
        {\textbf{itens do lote}, numerados\\(1 a $n$): as descrições\\
         a rotular};
    \draw[dashed, rounded corners=3pt] (-3.25,1.15) rectangle (3.25,-3.75);
    \node[rot, anchor=south west] at (-3.2,1.25) {uma chamada ao provedor};

    % ---------------- oráculo e resposta -----------------------------------
    \node[caixa, text width=28mm] (llm) at (5.6,-1.3)
        {oráculo LLM\\(provedor)};
    \node[resp, text width=48mm] (vet) at (10.2,-1.3)
        {\textbf{resposta}: vetor indexado,\\um objeto por item\\
         (categoria restrita às 621;\\expansão; justificativa)};

    \draw[seta] (3.25,-1.3) -- (llm.west);
    \draw[seta] (llm.east) -- (vet.west);

    % ---------------- nota do cache ----------------------------------------
    \node[rot, anchor=west, text width=128mm, align=left] at (-3.2,-4.7)
        {o prefixo repetido habilita o \textit{cache} de prefixo dos
         provedores; o esquema do lote não fixa o número de itens,
         preservando a cacheabilidade.};
\end{tikzpicture}
\end{document}
% GEOMETRIA CALIBRADA PARA CORPO 12 (ppginf/book 12pt, textwidth 16cm).
% Requer apenas arrows.meta e positioning (já em packages.tex). Legível em
% P&B: prefixo cacheável = cinza claro; resposta = cinza médio;
% contorno tracejado = a chamada.
% ============================================================================
\begin{tikzpicture}[
    x=1cm, y=1cm,
    caixa/.style={draw, rounded corners=2pt, align=center, text width=56mm,
                  minimum height=12mm, inner sep=3pt, font=\footnotesize},
    pref/.style={caixa, fill=gray!8},
    resp/.style={caixa, fill=gray!14},
    seta/.style={-{Stealth[length=2.6mm]}, thick},
    rot/.style={font=\scriptsize, align=center, inner sep=2pt}
]
    % ---------------- a chamada (contorno tracejado) -----------------------
    \node[pref] (prefixo) at (0,0)
        {\textbf{prefixo estático}, idêntico\\entre chamadas: contexto\\
         da tarefa $+$ esquema de saída};
    \node[caixa] (itens) at (0,-2.6)
        {\textbf{itens do lote}, numerados\\(1 a $n$): as descrições\\
         a rotular};
    \draw[dashed, rounded corners=3pt] (-3.25,1.15) rectangle (3.25,-3.75);
    \node[rot, anchor=south west] at (-3.2,1.25) {uma chamada ao provedor};

    % ---------------- oráculo e resposta -----------------------------------
    \node[caixa, text width=28mm] (llm) at (5.6,-1.3)
        {oráculo LLM\\(provedor)};
    \node[resp, text width=48mm] (vet) at (10.2,-1.3)
        {\textbf{resposta}: vetor indexado,\\um objeto por item\\
         (categoria restrita às 621;\\expansão; justificativa)};

    \draw[seta] (3.25,-1.3) -- (llm.west);
    \draw[seta] (llm.east) -- (vet.west);

    % ---------------- nota do cache ----------------------------------------
    \node[rot, anchor=west, text width=128mm, align=left] at (-3.2,-4.7)
        {o prefixo repetido habilita o \textit{cache} de prefixo dos
         provedores; o esquema do lote não fixa o número de itens,
         preservando a cacheabilidade.};
\end{tikzpicture}

% AUTOCONTIDO: sem \ref, sem códigos de experimento; o único número (621)
% é o do CategorySchema, definido na tese e usado no próprio apêndice.
% Uso standalone (pré-visualização):
%   \documentclass[tikz,border=6pt,12pt]{standalone}
%   \usetikzlibrary{arrows.meta, positioning}
%   \begin{document}% ============================================================================
% ESQUEMA (banca): anatomia de uma chamada de rotulagem em lote ao oráculo
% LLM — prefixo estático cacheável + itens numerados -> resposta como vetor
% indexado. Inserido no APÊNDICE E (a5-prompts/texto.tex) por ordem do
% autor (2026-08-23), dentro de figure:
%   % ============================================================================
% ESQUEMA (banca): anatomia de uma chamada de rotulagem em lote ao oráculo
% LLM — prefixo estático cacheável + itens numerados -> resposta como vetor
% indexado. Inserido no APÊNDICE E (a5-prompts/texto.tex) por ordem do
% autor (2026-08-23), dentro de figure:
%   % ============================================================================
% ESQUEMA (banca): anatomia de uma chamada de rotulagem em lote ao oráculo
% LLM — prefixo estático cacheável + itens numerados -> resposta como vetor
% indexado. Inserido no APÊNDICE E (a5-prompts/texto.tex) por ordem do
% autor (2026-08-23), dentro de figure:
%   \input{3-metodo/esquemas-propostos/esq-prompt-anatomia}
% AUTOCONTIDO: sem \ref, sem códigos de experimento; o único número (621)
% é o do CategorySchema, definido na tese e usado no próprio apêndice.
% Uso standalone (pré-visualização):
%   \documentclass[tikz,border=6pt,12pt]{standalone}
%   \usetikzlibrary{arrows.meta, positioning}
%   \begin{document}\input{esq-prompt-anatomia.tex}\end{document}
% GEOMETRIA CALIBRADA PARA CORPO 12 (ppginf/book 12pt, textwidth 16cm).
% Requer apenas arrows.meta e positioning (já em packages.tex). Legível em
% P&B: prefixo cacheável = cinza claro; resposta = cinza médio;
% contorno tracejado = a chamada.
% ============================================================================
\begin{tikzpicture}[
    x=1cm, y=1cm,
    caixa/.style={draw, rounded corners=2pt, align=center, text width=56mm,
                  minimum height=12mm, inner sep=3pt, font=\footnotesize},
    pref/.style={caixa, fill=gray!8},
    resp/.style={caixa, fill=gray!14},
    seta/.style={-{Stealth[length=2.6mm]}, thick},
    rot/.style={font=\scriptsize, align=center, inner sep=2pt}
]
    % ---------------- a chamada (contorno tracejado) -----------------------
    \node[pref] (prefixo) at (0,0)
        {\textbf{prefixo estático}, idêntico\\entre chamadas: contexto\\
         da tarefa $+$ esquema de saída};
    \node[caixa] (itens) at (0,-2.6)
        {\textbf{itens do lote}, numerados\\(1 a $n$): as descrições\\
         a rotular};
    \draw[dashed, rounded corners=3pt] (-3.25,1.15) rectangle (3.25,-3.75);
    \node[rot, anchor=south west] at (-3.2,1.25) {uma chamada ao provedor};

    % ---------------- oráculo e resposta -----------------------------------
    \node[caixa, text width=28mm] (llm) at (5.6,-1.3)
        {oráculo LLM\\(provedor)};
    \node[resp, text width=48mm] (vet) at (10.2,-1.3)
        {\textbf{resposta}: vetor indexado,\\um objeto por item\\
         (categoria restrita às 621;\\expansão; justificativa)};

    \draw[seta] (3.25,-1.3) -- (llm.west);
    \draw[seta] (llm.east) -- (vet.west);

    % ---------------- nota do cache ----------------------------------------
    \node[rot, anchor=west, text width=128mm, align=left] at (-3.2,-4.7)
        {o prefixo repetido habilita o \textit{cache} de prefixo dos
         provedores; o esquema do lote não fixa o número de itens,
         preservando a cacheabilidade.};
\end{tikzpicture}

% AUTOCONTIDO: sem \ref, sem códigos de experimento; o único número (621)
% é o do CategorySchema, definido na tese e usado no próprio apêndice.
% Uso standalone (pré-visualização):
%   \documentclass[tikz,border=6pt,12pt]{standalone}
%   \usetikzlibrary{arrows.meta, positioning}
%   \begin{document}% ============================================================================
% ESQUEMA (banca): anatomia de uma chamada de rotulagem em lote ao oráculo
% LLM — prefixo estático cacheável + itens numerados -> resposta como vetor
% indexado. Inserido no APÊNDICE E (a5-prompts/texto.tex) por ordem do
% autor (2026-08-23), dentro de figure:
%   \input{3-metodo/esquemas-propostos/esq-prompt-anatomia}
% AUTOCONTIDO: sem \ref, sem códigos de experimento; o único número (621)
% é o do CategorySchema, definido na tese e usado no próprio apêndice.
% Uso standalone (pré-visualização):
%   \documentclass[tikz,border=6pt,12pt]{standalone}
%   \usetikzlibrary{arrows.meta, positioning}
%   \begin{document}\input{esq-prompt-anatomia.tex}\end{document}
% GEOMETRIA CALIBRADA PARA CORPO 12 (ppginf/book 12pt, textwidth 16cm).
% Requer apenas arrows.meta e positioning (já em packages.tex). Legível em
% P&B: prefixo cacheável = cinza claro; resposta = cinza médio;
% contorno tracejado = a chamada.
% ============================================================================
\begin{tikzpicture}[
    x=1cm, y=1cm,
    caixa/.style={draw, rounded corners=2pt, align=center, text width=56mm,
                  minimum height=12mm, inner sep=3pt, font=\footnotesize},
    pref/.style={caixa, fill=gray!8},
    resp/.style={caixa, fill=gray!14},
    seta/.style={-{Stealth[length=2.6mm]}, thick},
    rot/.style={font=\scriptsize, align=center, inner sep=2pt}
]
    % ---------------- a chamada (contorno tracejado) -----------------------
    \node[pref] (prefixo) at (0,0)
        {\textbf{prefixo estático}, idêntico\\entre chamadas: contexto\\
         da tarefa $+$ esquema de saída};
    \node[caixa] (itens) at (0,-2.6)
        {\textbf{itens do lote}, numerados\\(1 a $n$): as descrições\\
         a rotular};
    \draw[dashed, rounded corners=3pt] (-3.25,1.15) rectangle (3.25,-3.75);
    \node[rot, anchor=south west] at (-3.2,1.25) {uma chamada ao provedor};

    % ---------------- oráculo e resposta -----------------------------------
    \node[caixa, text width=28mm] (llm) at (5.6,-1.3)
        {oráculo LLM\\(provedor)};
    \node[resp, text width=48mm] (vet) at (10.2,-1.3)
        {\textbf{resposta}: vetor indexado,\\um objeto por item\\
         (categoria restrita às 621;\\expansão; justificativa)};

    \draw[seta] (3.25,-1.3) -- (llm.west);
    \draw[seta] (llm.east) -- (vet.west);

    % ---------------- nota do cache ----------------------------------------
    \node[rot, anchor=west, text width=128mm, align=left] at (-3.2,-4.7)
        {o prefixo repetido habilita o \textit{cache} de prefixo dos
         provedores; o esquema do lote não fixa o número de itens,
         preservando a cacheabilidade.};
\end{tikzpicture}
\end{document}
% GEOMETRIA CALIBRADA PARA CORPO 12 (ppginf/book 12pt, textwidth 16cm).
% Requer apenas arrows.meta e positioning (já em packages.tex). Legível em
% P&B: prefixo cacheável = cinza claro; resposta = cinza médio;
% contorno tracejado = a chamada.
% ============================================================================
\begin{tikzpicture}[
    x=1cm, y=1cm,
    caixa/.style={draw, rounded corners=2pt, align=center, text width=56mm,
                  minimum height=12mm, inner sep=3pt, font=\footnotesize},
    pref/.style={caixa, fill=gray!8},
    resp/.style={caixa, fill=gray!14},
    seta/.style={-{Stealth[length=2.6mm]}, thick},
    rot/.style={font=\scriptsize, align=center, inner sep=2pt}
]
    % ---------------- a chamada (contorno tracejado) -----------------------
    \node[pref] (prefixo) at (0,0)
        {\textbf{prefixo estático}, idêntico\\entre chamadas: contexto\\
         da tarefa $+$ esquema de saída};
    \node[caixa] (itens) at (0,-2.6)
        {\textbf{itens do lote}, numerados\\(1 a $n$): as descrições\\
         a rotular};
    \draw[dashed, rounded corners=3pt] (-3.25,1.15) rectangle (3.25,-3.75);
    \node[rot, anchor=south west] at (-3.2,1.25) {uma chamada ao provedor};

    % ---------------- oráculo e resposta -----------------------------------
    \node[caixa, text width=28mm] (llm) at (5.6,-1.3)
        {oráculo LLM\\(provedor)};
    \node[resp, text width=48mm] (vet) at (10.2,-1.3)
        {\textbf{resposta}: vetor indexado,\\um objeto por item\\
         (categoria restrita às 621;\\expansão; justificativa)};

    \draw[seta] (3.25,-1.3) -- (llm.west);
    \draw[seta] (llm.east) -- (vet.west);

    % ---------------- nota do cache ----------------------------------------
    \node[rot, anchor=west, text width=128mm, align=left] at (-3.2,-4.7)
        {o prefixo repetido habilita o \textit{cache} de prefixo dos
         provedores; o esquema do lote não fixa o número de itens,
         preservando a cacheabilidade.};
\end{tikzpicture}

% AUTOCONTIDO: sem \ref, sem códigos de experimento; o único número (621)
% é o do CategorySchema, definido na tese e usado no próprio apêndice.
% Uso standalone (pré-visualização):
%   \documentclass[tikz,border=6pt,12pt]{standalone}
%   \usetikzlibrary{arrows.meta, positioning}
%   \begin{document}% ============================================================================
% ESQUEMA (banca): anatomia de uma chamada de rotulagem em lote ao oráculo
% LLM — prefixo estático cacheável + itens numerados -> resposta como vetor
% indexado. Inserido no APÊNDICE E (a5-prompts/texto.tex) por ordem do
% autor (2026-08-23), dentro de figure:
%   % ============================================================================
% ESQUEMA (banca): anatomia de uma chamada de rotulagem em lote ao oráculo
% LLM — prefixo estático cacheável + itens numerados -> resposta como vetor
% indexado. Inserido no APÊNDICE E (a5-prompts/texto.tex) por ordem do
% autor (2026-08-23), dentro de figure:
%   \input{3-metodo/esquemas-propostos/esq-prompt-anatomia}
% AUTOCONTIDO: sem \ref, sem códigos de experimento; o único número (621)
% é o do CategorySchema, definido na tese e usado no próprio apêndice.
% Uso standalone (pré-visualização):
%   \documentclass[tikz,border=6pt,12pt]{standalone}
%   \usetikzlibrary{arrows.meta, positioning}
%   \begin{document}\input{esq-prompt-anatomia.tex}\end{document}
% GEOMETRIA CALIBRADA PARA CORPO 12 (ppginf/book 12pt, textwidth 16cm).
% Requer apenas arrows.meta e positioning (já em packages.tex). Legível em
% P&B: prefixo cacheável = cinza claro; resposta = cinza médio;
% contorno tracejado = a chamada.
% ============================================================================
\begin{tikzpicture}[
    x=1cm, y=1cm,
    caixa/.style={draw, rounded corners=2pt, align=center, text width=56mm,
                  minimum height=12mm, inner sep=3pt, font=\footnotesize},
    pref/.style={caixa, fill=gray!8},
    resp/.style={caixa, fill=gray!14},
    seta/.style={-{Stealth[length=2.6mm]}, thick},
    rot/.style={font=\scriptsize, align=center, inner sep=2pt}
]
    % ---------------- a chamada (contorno tracejado) -----------------------
    \node[pref] (prefixo) at (0,0)
        {\textbf{prefixo estático}, idêntico\\entre chamadas: contexto\\
         da tarefa $+$ esquema de saída};
    \node[caixa] (itens) at (0,-2.6)
        {\textbf{itens do lote}, numerados\\(1 a $n$): as descrições\\
         a rotular};
    \draw[dashed, rounded corners=3pt] (-3.25,1.15) rectangle (3.25,-3.75);
    \node[rot, anchor=south west] at (-3.2,1.25) {uma chamada ao provedor};

    % ---------------- oráculo e resposta -----------------------------------
    \node[caixa, text width=28mm] (llm) at (5.6,-1.3)
        {oráculo LLM\\(provedor)};
    \node[resp, text width=48mm] (vet) at (10.2,-1.3)
        {\textbf{resposta}: vetor indexado,\\um objeto por item\\
         (categoria restrita às 621;\\expansão; justificativa)};

    \draw[seta] (3.25,-1.3) -- (llm.west);
    \draw[seta] (llm.east) -- (vet.west);

    % ---------------- nota do cache ----------------------------------------
    \node[rot, anchor=west, text width=128mm, align=left] at (-3.2,-4.7)
        {o prefixo repetido habilita o \textit{cache} de prefixo dos
         provedores; o esquema do lote não fixa o número de itens,
         preservando a cacheabilidade.};
\end{tikzpicture}

% AUTOCONTIDO: sem \ref, sem códigos de experimento; o único número (621)
% é o do CategorySchema, definido na tese e usado no próprio apêndice.
% Uso standalone (pré-visualização):
%   \documentclass[tikz,border=6pt,12pt]{standalone}
%   \usetikzlibrary{arrows.meta, positioning}
%   \begin{document}% ============================================================================
% ESQUEMA (banca): anatomia de uma chamada de rotulagem em lote ao oráculo
% LLM — prefixo estático cacheável + itens numerados -> resposta como vetor
% indexado. Inserido no APÊNDICE E (a5-prompts/texto.tex) por ordem do
% autor (2026-08-23), dentro de figure:
%   \input{3-metodo/esquemas-propostos/esq-prompt-anatomia}
% AUTOCONTIDO: sem \ref, sem códigos de experimento; o único número (621)
% é o do CategorySchema, definido na tese e usado no próprio apêndice.
% Uso standalone (pré-visualização):
%   \documentclass[tikz,border=6pt,12pt]{standalone}
%   \usetikzlibrary{arrows.meta, positioning}
%   \begin{document}\input{esq-prompt-anatomia.tex}\end{document}
% GEOMETRIA CALIBRADA PARA CORPO 12 (ppginf/book 12pt, textwidth 16cm).
% Requer apenas arrows.meta e positioning (já em packages.tex). Legível em
% P&B: prefixo cacheável = cinza claro; resposta = cinza médio;
% contorno tracejado = a chamada.
% ============================================================================
\begin{tikzpicture}[
    x=1cm, y=1cm,
    caixa/.style={draw, rounded corners=2pt, align=center, text width=56mm,
                  minimum height=12mm, inner sep=3pt, font=\footnotesize},
    pref/.style={caixa, fill=gray!8},
    resp/.style={caixa, fill=gray!14},
    seta/.style={-{Stealth[length=2.6mm]}, thick},
    rot/.style={font=\scriptsize, align=center, inner sep=2pt}
]
    % ---------------- a chamada (contorno tracejado) -----------------------
    \node[pref] (prefixo) at (0,0)
        {\textbf{prefixo estático}, idêntico\\entre chamadas: contexto\\
         da tarefa $+$ esquema de saída};
    \node[caixa] (itens) at (0,-2.6)
        {\textbf{itens do lote}, numerados\\(1 a $n$): as descrições\\
         a rotular};
    \draw[dashed, rounded corners=3pt] (-3.25,1.15) rectangle (3.25,-3.75);
    \node[rot, anchor=south west] at (-3.2,1.25) {uma chamada ao provedor};

    % ---------------- oráculo e resposta -----------------------------------
    \node[caixa, text width=28mm] (llm) at (5.6,-1.3)
        {oráculo LLM\\(provedor)};
    \node[resp, text width=48mm] (vet) at (10.2,-1.3)
        {\textbf{resposta}: vetor indexado,\\um objeto por item\\
         (categoria restrita às 621;\\expansão; justificativa)};

    \draw[seta] (3.25,-1.3) -- (llm.west);
    \draw[seta] (llm.east) -- (vet.west);

    % ---------------- nota do cache ----------------------------------------
    \node[rot, anchor=west, text width=128mm, align=left] at (-3.2,-4.7)
        {o prefixo repetido habilita o \textit{cache} de prefixo dos
         provedores; o esquema do lote não fixa o número de itens,
         preservando a cacheabilidade.};
\end{tikzpicture}
\end{document}
% GEOMETRIA CALIBRADA PARA CORPO 12 (ppginf/book 12pt, textwidth 16cm).
% Requer apenas arrows.meta e positioning (já em packages.tex). Legível em
% P&B: prefixo cacheável = cinza claro; resposta = cinza médio;
% contorno tracejado = a chamada.
% ============================================================================
\begin{tikzpicture}[
    x=1cm, y=1cm,
    caixa/.style={draw, rounded corners=2pt, align=center, text width=56mm,
                  minimum height=12mm, inner sep=3pt, font=\footnotesize},
    pref/.style={caixa, fill=gray!8},
    resp/.style={caixa, fill=gray!14},
    seta/.style={-{Stealth[length=2.6mm]}, thick},
    rot/.style={font=\scriptsize, align=center, inner sep=2pt}
]
    % ---------------- a chamada (contorno tracejado) -----------------------
    \node[pref] (prefixo) at (0,0)
        {\textbf{prefixo estático}, idêntico\\entre chamadas: contexto\\
         da tarefa $+$ esquema de saída};
    \node[caixa] (itens) at (0,-2.6)
        {\textbf{itens do lote}, numerados\\(1 a $n$): as descrições\\
         a rotular};
    \draw[dashed, rounded corners=3pt] (-3.25,1.15) rectangle (3.25,-3.75);
    \node[rot, anchor=south west] at (-3.2,1.25) {uma chamada ao provedor};

    % ---------------- oráculo e resposta -----------------------------------
    \node[caixa, text width=28mm] (llm) at (5.6,-1.3)
        {oráculo LLM\\(provedor)};
    \node[resp, text width=48mm] (vet) at (10.2,-1.3)
        {\textbf{resposta}: vetor indexado,\\um objeto por item\\
         (categoria restrita às 621;\\expansão; justificativa)};

    \draw[seta] (3.25,-1.3) -- (llm.west);
    \draw[seta] (llm.east) -- (vet.west);

    % ---------------- nota do cache ----------------------------------------
    \node[rot, anchor=west, text width=128mm, align=left] at (-3.2,-4.7)
        {o prefixo repetido habilita o \textit{cache} de prefixo dos
         provedores; o esquema do lote não fixa o número de itens,
         preservando a cacheabilidade.};
\end{tikzpicture}
\end{document}
% GEOMETRIA CALIBRADA PARA CORPO 12 (ppginf/book 12pt, textwidth 16cm).
% Requer apenas arrows.meta e positioning (já em packages.tex). Legível em
% P&B: prefixo cacheável = cinza claro; resposta = cinza médio;
% contorno tracejado = a chamada.
% ============================================================================
\begin{tikzpicture}[
    x=1cm, y=1cm,
    caixa/.style={draw, rounded corners=2pt, align=center, text width=56mm,
                  minimum height=12mm, inner sep=3pt, font=\footnotesize},
    pref/.style={caixa, fill=gray!8},
    resp/.style={caixa, fill=gray!14},
    seta/.style={-{Stealth[length=2.6mm]}, thick},
    rot/.style={font=\scriptsize, align=center, inner sep=2pt}
]
    % ---------------- a chamada (contorno tracejado) -----------------------
    \node[pref] (prefixo) at (0,0)
        {\textbf{prefixo estático}, idêntico\\entre chamadas: contexto\\
         da tarefa $+$ esquema de saída};
    \node[caixa] (itens) at (0,-2.6)
        {\textbf{itens do lote}, numerados\\(1 a $n$): as descrições\\
         a rotular};
    \draw[dashed, rounded corners=3pt] (-3.25,1.15) rectangle (3.25,-3.75);
    \node[rot, anchor=south west] at (-3.2,1.25) {uma chamada ao provedor};

    % ---------------- oráculo e resposta -----------------------------------
    \node[caixa, text width=28mm] (llm) at (5.6,-1.3)
        {oráculo LLM\\(provedor)};
    \node[resp, text width=48mm] (vet) at (10.2,-1.3)
        {\textbf{resposta}: vetor indexado,\\um objeto por item\\
         (categoria restrita às 621;\\expansão; justificativa)};

    \draw[seta] (3.25,-1.3) -- (llm.west);
    \draw[seta] (llm.east) -- (vet.west);

    % ---------------- nota do cache ----------------------------------------
    \node[rot, anchor=west, text width=128mm, align=left] at (-3.2,-4.7)
        {o prefixo repetido habilita o \textit{cache} de prefixo dos
         provedores; o esquema do lote não fixa o número de itens,
         preservando a cacheabilidade.};
\end{tikzpicture}
\end{document}
% GEOMETRIA CALIBRADA PARA CORPO 12 (ppginf/book 12pt, textwidth 16cm).
% Requer apenas arrows.meta e positioning (já em packages.tex). Legível em
% P&B: prefixo cacheável = cinza claro; resposta = cinza médio;
% contorno tracejado = a chamada.
% ============================================================================
\begin{tikzpicture}[
    x=1cm, y=1cm,
    caixa/.style={draw, rounded corners=2pt, align=center, text width=56mm,
                  minimum height=12mm, inner sep=3pt, font=\footnotesize},
    pref/.style={caixa, fill=gray!8},
    resp/.style={caixa, fill=gray!14},
    seta/.style={-{Stealth[length=2.6mm]}, thick},
    rot/.style={font=\scriptsize, align=center, inner sep=2pt}
]
    % ---------------- a chamada (contorno tracejado) -----------------------
    \node[pref] (prefixo) at (0,0)
        {\textbf{prefixo estático}, idêntico\\entre chamadas: contexto\\
         da tarefa $+$ esquema de saída};
    \node[caixa] (itens) at (0,-2.6)
        {\textbf{itens do lote}, numerados\\(1 a $n$): as descrições\\
         a rotular};
    \draw[dashed, rounded corners=3pt] (-3.25,1.15) rectangle (3.25,-3.75);
    \node[rot, anchor=south west] at (-3.2,1.25) {uma chamada ao provedor};

    % ---------------- oráculo e resposta -----------------------------------
    \node[caixa, text width=28mm] (llm) at (5.6,-1.3)
        {oráculo LLM\\(provedor)};
    \node[resp, text width=48mm] (vet) at (10.2,-1.3)
        {\textbf{resposta}: vetor indexado,\\um objeto por item\\
         (categoria restrita às 621;\\expansão; justificativa)};

    \draw[seta] (3.25,-1.3) -- (llm.west);
    \draw[seta] (llm.east) -- (vet.west);

    % ---------------- nota do cache ----------------------------------------
    \node[rot, anchor=west, text width=128mm, align=left] at (-3.2,-4.7)
        {o prefixo repetido habilita o \textit{cache} de prefixo dos
         provedores; o esquema do lote não fixa o número de itens,
         preservando a cacheabilidade.};
\end{tikzpicture}
